\section{Data 扩展：把互联网变成训练语料}

\subsection{能力决定数据}

如果你想让模型具备多语言、代码或数学能力，就必须问：数据从哪里来？怎么混配？怎么过滤？

\begin{figure}[H]
\centering
\begin{tikzpicture}[node distance=0.95cm, every node/.style={font=\small, align=center}]
\node[draw, rounded corners, fill=cyan!10, minimum width=4.0cm, minimum height=0.8cm] (goal) {desired capabilities};
\node[draw, rounded corners, fill=cyan!18, below=of goal, minimum width=4.0cm, minimum height=0.8cm] (src) {sources\\web / books / arXiv / code};
\node[draw, rounded corners, fill=cyan!26, below=of src, minimum width=4.0cm, minimum height=0.8cm] (proc) {processing\\transform / filter / dedup};
\node[draw, rounded corners, fill=cyan!34, below=of proc, minimum width=4.0cm, minimum height=0.8cm] (mix) {training mixture};
\draw[-{Latex[length=3mm]}, thick] (goal) -- (src);
\draw[-{Latex[length=3mm]}, thick] (src) -- (proc);
\draw[-{Latex[length=3mm]}, thick] (proc) -- (mix);
\end{tikzpicture}
\caption{数据工程的逻辑：目标能力决定来源，来源决定处理，处理决定混配}
\end{figure}

\begin{knowledgebox}{数据不是“抓到就行”}
\begin{itemize}
    \item 互联网上的数据是脏的、重复的、带偏差的。
    \item 格式不是文本而已，还包括 HTML、PDF、目录结构等。
    \item 你必须做转换、过滤和去重，才能让数据真的能训练。
\end{itemize}
\end{knowledgebox}

\subsection{evaluation / curation / processing}

课程把 data 模块拆成三层：先知道模型会什么，再决定要什么数据，再把数据处理成可训练样本。

\begin{figure}[H]
\centering
\begin{tikzpicture}[node distance=1.0cm, every node/.style={font=\small, align=center}]
\node[draw, rounded corners, fill=orange!12, minimum width=3.4cm, minimum height=0.8cm] (ev) {evaluation};
\node[draw, rounded corners, fill=orange!20, below=of ev, minimum width=3.4cm, minimum height=0.8cm] (cu) {curation};
\node[draw, rounded corners, fill=orange!28, below=of cu, minimum width=3.4cm, minimum height=0.8cm] (pr) {processing};
\node[draw, rounded corners, fill=orange!36, below=of pr, minimum width=3.4cm, minimum height=0.8cm] (tr) {training};
\draw[-{Latex[length=3mm]}, thick] (ev) -- (cu);
\draw[-{Latex[length=3mm]}, thick] (cu) -- (pr);
\draw[-{Latex[length=3mm]}, thick] (pr) -- (tr);
\end{tikzpicture}
\caption{data 模块的四步链路：评估、整理、处理、训练}
\end{figure}

\begin{table}[H]
\centering
\begin{tabular}{p{3cm}p{4cm}p{5.8cm}}
\toprule
\textbf{阶段} & \textbf{问题} & \textbf{常见方法} \\
\midrule
Evaluation & 模型会不会做某类任务 & perplexity、MMLU、instruction following \\
Curation & 哪些来源值得用 & web、books、code、papers、licensed data \\
Processing & 如何变成训练样本 & HTML/PDF to text、filter、dedup \\
\bottomrule
\end{tabular}
\caption{data 章节要训练你的判断顺序}
\end{table}

\begin{warningbox}{数据处理的第一原则}
先怀疑数据，再相信模型。
如果你不控制噪声、偏差和重复，模型会把这些东西原样学进去。
\end{warningbox}

